% 3.2 =============== Markov ===============

\section{马尔可夫过程与条件期望}    \label{sec-markov}

本节对马尔可夫过程理论进行讲解，并将其引入随机最优增长。
我们曾指出，处理连续动态规划的总体思路主要有离散化估计和平滑估计，例如VFI和Collocation，
但无论选取何种方法，\textbf{计算的主要难点之一}都在于期望（积分）的计算。
之所以先对马尔可夫过程进行学习，原因就在于它在两种方法中都发挥着重要作用：\textbf{它决定了条件期望的结构仅依赖当前状态。}

\begin{definition}
    \textbf{随机过程}（Stochastic Process）是某一\textbf{随机变量}（Random Variable）$z_t$构成的时间序列$\{z_t\}_{t=0}^{\infty}$。
    若$z_t$的\textbf{状态空间}$\mathbb{Z}$是有限或可数的，则称这一随机过程是离散取值的。
\end{definition}

\begin{definition}
    随机过程$\{z_t\}_{t=0}^{\infty}$被称为\textbf{马尔可夫过程}，若对任意$k\geq1$和任意$t$，成立：
    \begin{equation*}
        \operatorname{Pr}\left(z_{t+1}\mid z_t,z_{t-1},\cdots,z_{t-k}\right)=\operatorname{Pr}\left(z_{t+1}\mid z_t\right)
    \end{equation*}
    即条件于当前$z_t$，未来独立于历史$z^{t-1}$。如果上述概率不依赖于时间$t$，则称这一马尔可夫过程是\textbf{时间齐次}的。
\end{definition}

\begin{definition}
    \textbf{马尔科夫链}（Markov Chain）是离散取值的马尔可夫过程，即状态空间$\mathbb{Z}$可数或有限。
\end{definition}

我们将主要考虑时间齐次的马尔可夫链。为了记号的方便，同时不失一般性，假设随机变量$z_t$的取值空间
$\mathbb{Z}=\{\mathbbm{z}_1,\mathbbm{z}_2,\cdots,\mathbbm{z}_n\}$。
\footnote{通常我们用大写字母如$Z_t$表示随机变量，用小写字母如$z_t$表示样本的实现。这里是为了记号的方便以及与随机增长模型中一致，
我们使用小写的$z_t$表示随机变量}
通常，$n$维马尔可夫链可用三元符号$\left(\boldsymbol{\mathbbm{z}},\boldsymbol{\pi}_0, \boldsymbol{P}\right)$表示。
其中：

（1）\textbf{状态空间}$\boldsymbol{\mathbbm{z}}=\left(\mathbbm{z}_1,\mathbbm{z}_2,\cdots,\mathbbm{z}_n\right)^{\prime}$
是包含状态所有可能取值的$n$维列向量；

（2）\textbf{初始分布}$\boldsymbol{\pi}_0=\left(\pi_{01},\pi_{01},\cdots,\pi_{0n}\right)^{\prime}$是一个$n$维列向量，
其元素$\pi_{0i}=\operatorname{Pr}\left(s_0=s_i\right)$表示初始状态$z_0=\mathbbm{z}_i$，满足$\sum_{i=1}^{n}\pi_{0i}=1$。
在应用中，通常选定某一$i$，使$\pi_{0i}=1$，即从确定的初始状态$\mathbbm{z}_i$开始。

（3）\textbf{转移概率矩阵}$\boldsymbol{P}$具有如下形式：
\begin{equation*}
    \boldsymbol{P}=\left(\begin{array}{cccccc}
        p_{11} & p_{12} & \cdots & \vdots & \cdots & p_{1 n} \\
        p_{21} & & & \vdots & & \vdots \\
        \vdots & & & \vdots & & \vdots \\
        p_{i 1} & \cdots & \cdots & p_{i j} & \cdots & p_{i n} \\
        \vdots & & & \vdots & & \vdots \\
        p_{n 1} & \cdots & \cdots & \vdots & \cdots & p_{n n}
        \end{array}\right)
\end{equation*}
满足：
\begin{equation*}   
    \sum_{j=1}^{n}p_{ij}=1 \quad 0\leq p_{ij}\leq 1 \quad \forall i
\end{equation*}
这一矩阵也被称为\textbf{随机矩阵}（Stochastic Matrix）。矩阵中的代表元素$\left(i,j\right)$元$p_{ij}$是一个条件概率，表示：
\begin{equation*}
    p_{ij}=\operatorname{Pr}\left(z_{t+1}=\mathbbm{z}_j\mid z_t = \mathbbm{z}_i\right)
\end{equation*}


% 3.4.1 ===== Law of Motion for Pr =====
\subsection{转移路径与概率分布的运动方程}

马尔科夫链假设能够极大简化模型转移路径的计算。给定转移概率矩阵$\boldsymbol{P}$，
我们可以计算$m=1,2,3,\cdots$步内从状态$\mathbbm{z}_i$转移到状态$\mathbbm{z}_j$的概率。例如，对2步转移：
\begin{equation*}
    \begin{aligned}
        \operatorname{Pr}\left(z_{t+2}=\mathbbm{z}_j \mid z_t=\mathbbm{z}_i\right) 
            &= \sum_{k=1}^n \operatorname{Pr}\left(z_{t+2}=\mathbbm{z}_j \mid z_{t+1}=\mathbbm{z}_k\right) 
                    \operatorname{Pr}\left(z_{t+1}=\mathbbm{z}_k \mid z_t=\mathbbm{z}_i\right) \\
            &= \sum_{k=1}^n p_{k j} p_{i k} = \sum_{k=1}^n p_{i k} p_{k j} \\
            &= p^{\left(2\right)}_{i j} 
    \end{aligned}
\end{equation*}
其中$p^{\left(2\right)}_{i j}$表示两步转移概率矩阵$\boldsymbol{P}^2=\boldsymbol{P}\boldsymbol{P}$的$(i,j)$元。
容易验证若$\boldsymbol{P}$是一个转移概率矩阵，则$\boldsymbol{P}$的$m$次幂$\boldsymbol{P}^{m}$也是一个转移概率矩阵。
一般而言，$\boldsymbol{P}^{m}$的$(i,j)$元不一定等于$\boldsymbol{P}$的$(i,j)$元的$m$次方，
即$p^{\left(m\right)}_{i j}\neq \left(p_{i j}\right)^m$。
进一步归纳后不难得到：
\begin{equation}
    \operatorname{Pr}\left(z_{t+m}= \mathbbm{z}_j\mid z_t=\mathbbm{z}_i\right) = p^{\left(m\right)}_{i j}
\end{equation}
其中右侧表示$m$步转移概率矩阵$\boldsymbol{P}^m$的$(i,j)$元。
更一般化的，我们有如下查普曼-柯尔莫哥洛夫方程（Chapman-Kolmogorov Equation）：
\begin{equation}
    \begin{aligned}
        \boldsymbol{P}^{m+k} &= \boldsymbol{P}^k\boldsymbol{P}^m \\
        & \Leftrightarrow \\
        p^{(m+k)}_{i j} &= \sum_{r=1}^{n}p^{(k)}_{i r}p^{(m)}_{r j}
    \end{aligned}
\end{equation}

令$n$维列向量$\boldsymbol{\pi}_t=\left(\pi_{t,1},\cdots,\pi_{t,n}\right)^{\prime}$
表示随机变量$z_t$其状态空间$\mathbb{S}$上的\textbf{无条件概率分布}，其中的元素：
\begin{equation*}
    \pi_{t,i} = \operatorname{Pr}\left(z_t = \mathbbm{z}_i\right), \quad i=1,\cdots n
\end{equation*}
由转移概率可得上述概率分布随时间的变化：
\begin{equation*}
    \begin{aligned}
        \pi_{1, i} &= \sum_{j=1}^n \operatorname{Pr}\left[z_1=\mathbbm{z}_i \mid z_0=\mathbbm{z}_j\right] 
                        \operatorname{Pr}\left[z_0=\mathbbm{z}_j\right] 
                    =  \sum_{j=1}^n \operatorname{Pr}\left[z_1=\mathbbm{z}_i \mid z_0=\mathbbm{z}_j\right]\pi_{0,j}  \\
        \vdots & \\
        \pi_{t, i} &= \sum_{j=1}^n \operatorname{Pr}\left[z_t=\mathbbm{z}_i \mid z_{t-1}=\mathbbm{z}_j\right] 
                        \operatorname{Pr}\left[z_{t-1}=\mathbbm{z}_j\right] 
                    =  \sum_{j=1}^n \operatorname{Pr}\left[z_t=\mathbbm{z}_i \mid z_{t-1}=\mathbbm{z}_j\right]\pi_{t-1,j}  \\
        \pi_{t+1, i} &= \sum_{j=1}^n \operatorname{Pr}\left[z_{t+1}=\mathbbm{z}_i \mid z_t=\mathbbm{z}_j\right] 
                        \operatorname{Pr}\left[z_t=\mathbbm{z}_j\right]
                    =  \sum_{j=1}^n \operatorname{Pr}\left[z_{t+1}=\mathbbm{z}_i \mid z_t=\mathbbm{z}_j\right]\pi_{t,j}
        \end{aligned}
\end{equation*}
以矩阵形式表示，可得\textbf{概率分布的运动方程}（Law of Motion）为：
\begin{equation*}
    \begin{aligned}
        \boldsymbol{\pi}_1^{\prime} &= \boldsymbol{\pi}_0^{\prime}\boldsymbol{P} \\
        \boldsymbol{\pi}_2^{\prime} &= \boldsymbol{\pi}_1^{\prime}\boldsymbol{P} 
            = \boldsymbol{\pi}_0^{\prime}\boldsymbol{P}^2 \\
        \vdots & \\
        \boldsymbol{\pi}_t^{\prime} &= \boldsymbol{\pi}_{t-1}^{\prime}\boldsymbol{P} 
            = \boldsymbol{\pi}_0^{\prime}\boldsymbol{P}^t
    \end{aligned}
\end{equation*}
从而有：
\begin{equation}    \label{eq-markov-pr-motion}
    \boldsymbol{\pi}_{t+1} = \boldsymbol{P}^{\prime}\boldsymbol{\pi}_t 
        = \left(\boldsymbol{P}^{\prime}\right)^{t+1}\boldsymbol{\pi}_0
\end{equation}


% 3.4.2 ===== Stationary Distribution =====
\subsection{平稳分布}

\begin{definition}
    \textbf{稳态/平稳分布}（Stationary Distribution）是指满足如下条件的一个无条件概率分布：
    \begin{equation*}
        \boldsymbol{\pi}_{t+1} = \boldsymbol{\pi}_t
    \end{equation*}
\end{definition}
根据概率分布的运动方程\ref{eq-markov-pr-motion}，平稳分布必须满足：
\begin{equation}    \label{eq-markov-stationary}
    \boldsymbol{\pi} = \boldsymbol{P}^{\prime}\boldsymbol{\pi}
\end{equation}
等价于：
\begin{equation}
    \left(\boldsymbol{I}-\boldsymbol{P}^{\prime}\right)\boldsymbol{\pi} = \boldsymbol{0}
\end{equation}
这表明：\textbf{平稳分布$\boldsymbol{\pi}$是矩阵$\boldsymbol{P}^{\prime}$对应特征值为1的特征向量}（正则化使$\sum_{i=1}^{n}\pi_i=1$）。

\begin{definition}
    $\boldsymbol{P}, \boldsymbol{\pi}$被称为\textbf{平稳马尔科夫链}，若分布$\boldsymbol{\pi}$满足\ref{eq-markov-stationary}式。
\end{definition}

事实上，$\boldsymbol{P}$是随机矩阵（元素非负并且$\sum_{1}^{n}p_{ij}=1$）保证了$\boldsymbol{P}$至少有一个特征根为1，
进而至少有一个特征向量$\boldsymbol{\pi}$是平稳分布。但是，平稳分布可能不唯一，因为可能有重复为1的特征根。
确保平稳分布唯一的一个充分条件是：$0<p_{i j}<1, ~ \forall i, j$。

\begin{example}
    考虑一个两状态马尔科夫链，转移矩阵为：
    \begin{equation*}
        \boldsymbol{P} = 
        \begin{bmatrix}
            1-p & p \\
            q & 1-q
        \end{bmatrix}
    \end{equation*}
    平稳分布$\boldsymbol{\pi}=\left(\pi^1, \pi^2\right)^{\prime}$是如下方程的解：
    \begin{equation*}
        \left[
            \begin{pmatrix}
                1 & 0 \\
                 0 & 1
            \end{pmatrix}
            -
            \begin{pmatrix}
                1-p & q \\
                p & 1-q
            \end{pmatrix}
        \right]
        \begin{pmatrix}
            \pi^1 \\
            \pi^2
        \end{pmatrix}
        =
        \begin{pmatrix}
            0 \\
            0
        \end{pmatrix}
    \end{equation*}
    可以解得：
    \begin{equation*}
        \begin{pmatrix}
            \pi^1 \\
            \pi^2
        \end{pmatrix}
        =
        \begin{pmatrix}
            \frac{q}{p+q} \\
            \frac{p}{p+q}
        \end{pmatrix}
    \end{equation*}
    不难看出，当$q\rightarrow 0$，状态2会吸收状态1。
\end{example}

有时，我们也会考虑连续支撑上的马尔可夫过程，此时原有概率分布的概念相应变为了\textbf{概率密度}。
假设$z_t$的\textbf{概率密度}为$\pi_t(z)$，根据转移方程有：
\begin{equation*}
    \pi_{t+1}\left(z^{\prime}\right)=\int p\left(z^{\prime} \mid z\right) \pi_t(z) dz
\end{equation*}
其中，$p\left(z^{\prime}\mid z\right)$是条件于$z_t=z$下$z_{t+1}=z^{\prime}$的\textbf{概率密度}。
\textbf{平稳的概率密度函数$\pi(z)$}满足如下不动点条件：
\begin{equation*}
    \pi(z^{\prime}) = \int  p\left(z^{\prime} \mid z\right) \pi(z) dz
\end{equation*}

\begin{example}
    假设$\{z_t\}$是一个线性高斯AR(1)过程：
    \begin{equation*}
        z_{t+1}=(1-\rho) \mu+\rho z_t+\sigma \varepsilon_{t+1}, \quad 
            \varepsilon_{t+1} \sim \operatorname{iid} ~ N(0,1)
    \end{equation*}
    那么：
    \begin{equation*}
        p\left(z^{\prime} \mid z\right)=\frac{1}{\sigma} \phi\left(\frac{z^{\prime}-(1-\rho) \mu-\rho z}{\sigma}\right)
    \end{equation*}
    其中，$\phi(\varepsilon)$是标准正态分布的概率密度函数：
    \begin{equation*}
        \phi(\varepsilon) \equiv \frac{1}{\sqrt{2 \pi}} e^{-\varepsilon^2 / 2}
    \end{equation*}
    若$\left|\rho\right|<1$，那么存在唯一的平稳概率密度：
    \begin{equation*}
        \pi(z) = \frac{\sqrt{1-\rho^2}}{\sigma} \phi\left(\sqrt{1-\rho^2}\frac{(z-\mu)}{\sigma}\right)
    \end{equation*}
\end{example}


% 3.4.3 ===== 马尔可夫假设下的随机最优增长 =====
\subsection{马尔可夫假设下的随机最优增长}

在马尔可夫假设下，$z_t$服从一个一阶马尔可夫过程，条件概率（密度）为$\pi(z^{\prime}\mid z)$，这使得未来不依赖于众多可能的历史
而只依赖于当前，进而我们可以写出递归形式的随机最优增长：
\begin{equation*}
    v(k, z)=\max _{k^{\prime}}
        \left[u\left(zf(k)+(1-\delta)k-k^{\prime}\right)
        +\beta \int v\left(k^{\prime}, z^{\prime}\right) \pi\left(z^{\prime} \mid z\right) d z^{\prime}\right]
\end{equation*}
不难得到欧拉泛函为：
\begin{equation*}
    u^{\prime}\left(z f(k)-k^{\prime}\right)=
        \beta \int u^{\prime}\left(z^{\prime} f\left(k^{\prime}\right) - 
            k^{\prime \prime}\right) z^{\prime} f^{\prime}\left(k^{\prime}\right) 
                \pi\left(z^{\prime} \mid z\right) d z^{\prime}
\end{equation*}
我们指出它隐式地定义了动态规划问题的策略函数$k^{\prime}=g(k,z)$。形式上，我们可以写出：
\begin{equation*}
    k_{t+1} = g(k_t,z_t)
\end{equation*}
现在，只从模型性质分析的角度出发，我们面临的一个问题是：$\{k_t\}$未必收敛至某一稳态，那么\textbf{$k$的分布是怎样的？}

假设策略函数具有乘积形式，即$k_{t+1}=z_tg(k_t)$，并且$z_t$随时间是独立同分布的，因此可记累积分布函数为：
\begin{equation*}
    H(z)\equiv \operatorname{Pr}[z_t\leq z]
\end{equation*}
接下来考虑$t$时刻资本$k$的累计分布函数：
\begin{equation*}
    \Psi_t(k) \equiv \operatorname{Pr}[k_t\leq k]
\end{equation*}
例如，对$t=1$：
\begin{equation*}
    \begin{aligned}
        \Psi_1(k) & =\operatorname{Prob}\left[k_1 \leq k\right] \\
            &= \operatorname{Prob}\left[z_0 g\left(k_0\right) \leq k\right]
                =\operatorname{Prob}\left[z_0 \leq \frac{k}{g\left(k_0\right)}\right] \\
            &= H\left(\frac{k}{g\left(k_0\right)}\right)
    \end{aligned}
\end{equation*}
记$P(k^{\prime}\mid k)$为条件分布：
\begin{equation*}
    P\left(k^{\prime} \mid k\right) \equiv \operatorname{Prob}\left[k_{t+1} \leq k^{\prime} \mid k_t=k\right]
\end{equation*}
在独立同分布的假设下，可以得到：
\begin{equation*}
    P(k^{\prime} \mid k) = H\left(\frac{k^{\prime}}{g(k)}\right)
\end{equation*}
资本$k$的累计分布函数$\Psi(\cdot)$的运动方程为：
\begin{equation*}
    \Psi_{t+1}\left(k^{\prime}\right)=\int P\left(k^{\prime} \mid k\right) d \Psi_t(k)
\end{equation*}
用概率密度函数$\psi(\cdot)$表示这一运动方程，则有：
\begin{equation*}
    \psi_{t+1}\left(k^{\prime}\right)=\int p\left(k^{\prime} \mid k\right) \psi_t(k) dk
\end{equation*}
平稳概率密度是这一运动方程的不动点，即：
\begin{equation*}
    \psi\left(k^{\prime}\right)=\int p\left(k^{\prime} \mid k\right) \psi(k) dk
\end{equation*}
更一般的来说，由策略函数$k^{\prime}=g(k,z)$和外生冲击的条件密度函数$\pi\left(z^{\prime}\mid z\right)$，
可以得到状态$(k,z)$的联合分布（Joint Distribution）的运动方程，我们希望找到其不动点。

马尔可夫性的另一个好处是简化了模型中条件期望的计算，因为它使得下期状态只依赖于当前而不依赖于历史。
例如，假设随机过程$\{z_t\}_{t=0}^{\infty}$是一个马尔科夫链，
取值空间为$\boldsymbol{\mathbbm{z}}=\left(\mathbbm{z}_1,\mathbbm{z}_2,\cdots,\mathbbm{z}_n\right)^{\prime}$。
那么给定$z_t=\mathbbm{z}_i$，$z_{t+1}$的条件期望为：
\begin{equation}
    \mathbb{E}\left(z_{t+1} \mid z_t=\mathbbm{z}_i\right) = \boldsymbol{P}_i\boldsymbol{\mathbbm{z}}
\end{equation}
其中，$\boldsymbol{P}_i$为转移概率矩阵$\boldsymbol{P}$的第$i$行（即条件于当前状态为$\mathbbm{z}_i$）。
对于随机最优增长模型，假设$z_t$服从一个$n$维马尔可夫链，取值空间同上。
给定$k=k_m,~z=z_i$，对值函数$v\left(k_m,z_i\right)$，利用转移概率可以直接写出如下贝尔曼方程：
\begin{equation*}
    \begin{aligned}
        v\left(k_m,z_i\right) &= \max_{k^{\prime}}u\left[zf(k_m)+(1-\delta)k_m-k^{\prime}\right] + 
            \beta\mathbbm{E}\left[v\left(k^{\prime},z^{\prime}\right)\mid z=z_i\right] \\
        &= \max_{k^{\prime}}u\left[zf(k_m)+(1-\delta)k_m-k^{\prime}\right] +
            \beta\boldsymbol{P}_i v\left(k^{\prime},\mathbbm{z}\right)
    \end{aligned}
\end{equation*}
其中$v\left(k^{\prime},\mathbbm{z}\right)$是一个$n$维列向量，每行对应下期$z^{\prime}$在空间$\mathbbm{z}$上可能的取值
$\left(z_1,\cdots,z_n\right)$，$\boldsymbol{P}_i$则是转移矩阵的第$i$行，它表示条件于当前TFP状态为$z_i$。

我们可以将其写为更为紧凑的矩阵形式。记$v\left(k_m,\mathbbm{z}\right)$是$n$维值函数，每行对应TFP的一种可能状态下的最优值。
令$\vec{\boldsymbol{1}}$是元素全部维1的$n$维列向量，$\boldsymbol{z}^{\prime}$表示下期的$n$维列向量，从而有：
\begin{equation*}
    \begin{aligned}
        v\left(k_m,\mathbbm{z}\right) 
        &= \max_{k^{\prime}}u\left[\mathbbm{z}f(k_m)+\vec{\boldsymbol{1}}(1-\delta)k_m-\vec{\boldsymbol{1}}k^{\prime}\right] + 
            \beta\mathbbm{E}\left[v\left(k^{\prime},\boldsymbol{z}^{\prime}\right)\mid \mathbbm{z}\right] \\
        &= \max_{k^{\prime}}u\left[\mathbbm{z}f(k_m)+\vec{\boldsymbol{1}}(1-\delta)k_m-\vec{\boldsymbol{1}}k^{\prime}\right] +
            \beta\boldsymbol{P} v\left(k^{\prime},\mathbbm{z}\right)
    \end{aligned}
\end{equation*}


